\begin{notechapter}
\chapter{Affine Spaces}
\label{ch:affine-spaces}
\chaptersubtitle{%
  Points, displacements, and what a parameter update is allowed to do}

\begin{chapterintro}
The update rule of gradient descent,
$\theta\leftarrow\theta - \eta\nabla L(\theta)$, adds two objects of
different kinds. The parameter $\theta$ is a \emph{point}:
a state of the model. The term $-\eta\nabla L(\theta)$ is a \emph{displacement}:
a change of state. A computer stores both as arrays of $n$ numbers, so the
difference is invisible in code, and most calculus books erase it on the
first page by declaring that both are elements of $\R^n$. This chapter keeps
them apart. The structure in which points and displacements live together
is called an affine space, and everything else in the book is built on it.
You will learn
\begin{itemize}[nosep]
  \item the definition of an affine space and the rules of calculation that
        follow from it: point $+$ vector is a point, point $-$ point is a
        vector, and nothing else is allowed;
  \item what choosing an origin does, and why coordinates of points and of
        vectors react differently when the origin moves;
  \item which combinations of points are meaningful (weights summing to $1$
        give a point, weights summing to $0$ give a vector);
  \item affine subspaces, convexity, and affine maps $x\mapsto Ax + b$, where
        $A$ belongs to the map and $b$ to the choice of origins;
  \item the \emph{origin test}, which detects formulas that secretly depend
        on the origin, such as weight decay.
\end{itemize}
\prerequisites{vector spaces, linear maps, bases and matrices, at the level
of a first course in linear algebra. \Cref{ch:duality} revisits them. All
vector spaces in this chapter are real; from
\cref{affine-spaces:sec:coordinates} on they are finite-dimensional.}
\end{chapterintro}

\section{Two kinds of data}\label{affine-spaces:sec:two-kinds}

Some quantities answer the question ``where?'' and others the question
``how far, and in which direction?''.

\begin{itemize}[nosep]
  \item \emph{Dates and durations.} ``14 July, noon'' is a date; ``three
        hours'' is a duration. A date plus a duration is a date, and the
        difference of two dates is a duration. Two durations can be added.
        Two dates cannot: ``14 July plus 3 August'' means nothing.
  \item \emph{Temperatures and temperature differences.} A temperature of
        $20^\circ$C is a level on a scale whose zero is conventional.
        ``$5$ degrees warmer'' is a difference. Twice ``$5$ degrees warmer''
        is ``$10$ degrees warmer'' on every scale with the same unit, but
        twice $20^\circ$C is $40^\circ$C, whereas the same computation in
        kelvins gives $586.3$~K, that is, $313.15^\circ$C. Doubling a
        temperature level depends on where the zero was put.
  \item \emph{Positions and displacements.} A place on a map is a position;
        ``$2$ km to the north-east'' is a displacement, and it can be
        applied at any position.
  \item \emph{Parameters and updates.} The weights of a model after $k$
        steps of training form a position in parameter space. The update
        applied at step $k$ is a displacement.
\end{itemize}

Programming languages that take types seriously enforce the distinction.
In Python, the difference of two \texttt{datetime} objects is a
\texttt{timedelta}, a \texttt{datetime} plus a \texttt{timedelta} is a
\texttt{datetime}, and the sum of two \texttt{datetime} objects raises a
\texttt{TypeError}. Tensors have no such protection: \texttt{x + y} is
accepted whatever \texttt{x} and \texttt{y} stand for. The mathematics has
to do the type checking, and the type system it uses is the definition of
the next section.

\begin{figure}[ht]
\centering
\begin{tikzpicture}[>=Stealth, scale=1.1]
  \fill (0,0) circle (1.6pt) node[below left] {$a$};
  \fill (2.4,1.0) circle (1.6pt) node[above right] {$b = a + v$};
  \draw[->, very thick, bookdef] (0,0) -- (2.4,1.0)
    node[midway, above left] {$v = b - a$};
  \fill (3.2,-0.9) circle (1.6pt) node[below left] {$c$};
  \fill (5.6,0.1) circle (1.6pt) node[above right] {$c + v$};
  \draw[->, very thick, bookdef] (3.2,-0.9) -- (5.6,0.1)
    node[midway, below right] {$v$};
\end{tikzpicture}
\caption{Points (black dots) and a displacement (green arrow). The
displacement $v$ that leads from $a$ to $b$ is not attached to $a$: applied
at another point $c$, the same $v$ leads to $c + v$.}
\label{affine-spaces:fig:points-vectors}
\end{figure}

\section{The definition}\label{affine-spaces:sec:definition}

\begin{definition}[Affine space]\label{affine-spaces:def:affine}
An \emph{affine space} is a non-empty set $\Aff$, whose elements are called
\emph{points}, together with a real vector space $V$, whose elements are
called \emph{displacements} or simply \emph{vectors}, and an operation
$\Aff\times V\to\Aff$, written $(a, v)\mapsto a + v$, such that for all
$a, b\in\Aff$ and $u, v\in V$:
\begin{enumerate}[nosep]
  \item $a + 0 = a$;
  \item $(a + u) + v = a + (u + v)$;
  \item there is exactly one $v\in V$ such that $a + v = b$; it is written
        $b - a$.
\end{enumerate}
The space $V$ is the \emph{direction} of $\Aff$, and the \emph{dimension} of
$\Aff$ is the dimension of $V$.
\end{definition}

\begin{intuition}
The vector space $V$ acts on the set $\Aff$ by moving its points. Axioms 1
and 2 say that the zero displacement moves nothing and that two successive
displacements amount to their sum. Axiom 3 says that from any point, every
other point is reached by one displacement and only one. The symbol $+$ is
used with two meanings, in $a + v$ and in $u + v$; axiom 2 is what makes
this harmless.
\end{intuition}

The definition provides two operations between the two sets,
\[
  \text{point} + \text{vector} = \text{point},
  \qquad
  \text{point} - \text{point} = \text{vector},
\]
and nothing else. It does not define the sum of two points, the product of
a point by a number, or a distinguished point that would deserve to be
called zero. All the usual rules that mix the two operations follow from
the three axioms.

\begin{proposition}[Rules of calculation]\label{affine-spaces:prop:rules}
For all $a, b, c\in\Aff$ and $u, v\in V$:
\begin{enumerate}[nosep]
  \item if $a + u = a + v$ then $u = v$; in particular $a + v = a$ only for
        $v = 0$, and $a - a = 0$;
  \item \emph{(Chasles' relation)} $(b - a) + (c - b) = c - a$;
  \item $a - b = -(b - a)$;
  \item $(a + u) - (b + v) = (a - b) + u - v$;
  \item if $a + u = b + u$ then $a = b$.
\end{enumerate}
\end{proposition}

\begin{proof}
(1) Put $b = a + u = a + v$. Both $u$ and $v$ are displacements from $a$ to
$b$, so $u = v$ by the uniqueness in axiom 3. Since $a + 0 = a$ by axiom 1,
the only $v$ with $a + v = a$ is $0$, which is the meaning of $a - a = 0$.

(2) By axiom 2 and the definition of $b - a$ and $c - b$,
\[
  a + \bigl((b - a) + (c - b)\bigr) = \bigl(a + (b - a)\bigr) + (c - b)
  = b + (c - b) = c .
\]
So $(b - a) + (c - b)$ is a displacement from $a$ to $c$, and by uniqueness
it is $c - a$.

(3) Take $c = a$ in (2): $(b - a) + (a - b) = a - a = 0$.

(4) Using axiom 2 and the commutativity of addition in $V$,
\[
  (b + v) + \bigl((a - b) + u - v\bigr) = b + \bigl((a - b) + u\bigr)
  = \bigl(b + (a - b)\bigr) + u = a + u .
\]
So $(a - b) + u - v$ is the displacement from $b + v$ to $a + u$.

(5) Add $-u$ to both sides: $a = (a + u) + (-u) = (b + u) + (-u) = b$, by
axioms 2 and 1.
\end{proof}

\begin{example}[A vector space is an affine space]
\label{affine-spaces:ex:vector-space}
Every vector space $V$ is an affine space with direction $V$ itself: the
points are the elements of $V$, and $a + v$ is the addition of $V$. Then
$b - a$ is the usual difference. This is how $\R^n$ is used in most of
calculus. In this example the set of points happens to contain a special
element, $0$, and to carry extra operations (sums and multiples of points).
The affine structure does not see them.
\end{example}

\begin{example}[Solutions of a linear system]\label{affine-spaces:ex:system}
Let $A$ be an $m\times n$ matrix and $b\in\R^m$, and suppose that the set
$S = \{x\in\R^n : Ax = b\}$ is not empty. If $b\neq0$, $S$ is not a vector
space: it does not contain $0$, and the sum of two solutions is not a
solution, since $A(x + y) = 2b$. But the difference of two solutions lies
in $\ker A$, and a solution plus an element of $\ker A$ is a solution. So
$S$ is an affine space with direction $\ker A$, and it has no natural
origin: no solution is more special than another.
\end{example}

\begin{example}[Probability vectors]\label{affine-spaces:ex:simplex}
The set $H = \{p\in\R^n : p_1 + \dots + p_n = 1\}$ is an affine space with
direction $H_0 = \{v\in\R^n : v_1 + \dots + v_n = 0\}$. Probability
distributions on $n$ outcomes are the points of $H$ with non-negative
coordinates. The sum of two probability vectors is not a probability
vector, but their difference is an element of $H_0$. The familiar quantity
$p - y$ (predicted probabilities minus one-hot label), which appears when
the cross-entropy is differentiated with respect to the logits, is such a
difference: its coordinates sum to zero.
\end{example}

\begin{pitfall}[Angles are not points of an affine space]
Angles, hours on a clock face and other periodic quantities can be moved by
a displacement ($350^\circ + 20^\circ = 10^\circ$), and axioms 1 and 2 hold.
Axiom 3 fails: from $350^\circ$ to $10^\circ$ the displacement can be
$+20^\circ$, $-340^\circ$, $+380^\circ$, and so on. Operations that rely on
uniqueness fail with it. The ``midpoint'' of $350^\circ$ and $10^\circ$
computed as $(350 + 10)/2 = 180^\circ$ points in the direction opposite to
the two angles, whereas $(-10 + 10)/2 = 0^\circ$ lies between them
(\cref{affine-spaces:exr:7}). The circle is the first example of a space
that looks affine near each point but not globally;
\cref{ch:tangent-bundles} returns to it.
\end{pitfall}

\begin{remark}[An equivalent definition]
Some authors start from the map $\Aff\times\Aff\to V$, $(a, b)\mapsto b - a$,
and ask for two properties: Chasles' relation, and that for every $a$ the
map $b\mapsto b - a$ is a bijection from $\Aff$ onto $V$. The operation
$a + v$ is then \emph{defined} as the unique $b$ with $b - a = v$. The two
definitions describe the same objects; the one adopted here puts forward
the operation that an update rule uses.
\end{remark}

\section{Origins and coordinates}\label{affine-spaces:sec:coordinates}

From now on $\Aff$ has finite dimension $n$.

\begin{definition}[Vectorialization, frame, coordinates]
\label{affine-spaces:def:frame}
Let $o\in\Aff$. The map
\[
  \Theta_o:\Aff\to V,\qquad a\mapsto a - o
\]
is the \emph{vectorialization} of $\Aff$ at $o$, and $o$ is then called an
\emph{origin}. An \emph{affine frame} of $\Aff$ is a pair
$(o; e_1,\dots,e_n)$ made of a point $o$ and a basis $(e_1,\dots,e_n)$ of
$V$. The \emph{coordinates} of a point $a$ in this frame are the numbers
$x_1,\dots,x_n$ such that $a = o + \sum_jx_je_j$, that is, the components
of the vector $a - o$ in the basis. The coordinates of a vector $v$ in the
frame are its components in the basis; the origin plays no role in them.
\end{definition}

By axiom 3, $\Theta_o$ is a bijection, with inverse $v\mapsto o + v$. Once
an origin is chosen, points can therefore be \emph{represented} by vectors,
and once a basis is chosen too, by elements of $\R^n$. This is what every
program does. The representation depends on the choices, and points and
vectors do not react in the same way when the origin changes.

\begin{proposition}[Moving the origin]\label{affine-spaces:prop:origin-change}
Let $o' = o + w$ be another origin, with the same basis. If a point $a$ has
coordinates $x$ in the frame $(o; e_1,\dots,e_n)$ and $x'$ in the frame
$(o'; e_1,\dots,e_n)$, and $w$ has components $\omega$, then
\[
  x' = x - \omega .
\]
The coordinates of a vector are the same in both frames.
\end{proposition}

\begin{proof}
By rule (4) of \cref{affine-spaces:prop:rules},
$a - o' = (a + 0) - (o + w) = (a - o) - w$. Taking components in the basis
gives $x' = x - \omega$. The coordinates of a vector are its components in
the basis, which has not changed.
\end{proof}

\begin{keyidea}
When the origin moves by $w$, the coordinates of every point change by
$-w$ and the coordinates of every vector stay the same. A formula written
in coordinates describes something geometric only if its result reacts
accordingly: like a point if it is meant to be a point, like a vector if
it is meant to be a vector, and not at all if it is meant to be a number.
\end{keyidea}

The temperature scales of \cref{affine-spaces:sec:two-kinds} are frames of a
one-dimensional affine space. Celsius and Kelvin share the basis vector
(one degree) and differ by the origin, with $\omega = -273.15$ when passing
from Celsius to Kelvin coordinates. Celsius and Fahrenheit differ by both
the origin and the basis vector.

\section{Affine combinations}\label{affine-spaces:sec:combinations}

Adding two points is not defined. Yet the midpoint of two points and the
mean of a cloud of points are perfectly geometric. The following
proposition explains which weighted sums of points make sense.

\begin{proposition}[Affine combinations]\label{affine-spaces:prop:affine-comb}
Let $a_1,\dots,a_k\in\Aff$, $\lambda_1,\dots,\lambda_k\in\R$ and $o\in\Aff$.
\begin{enumerate}[nosep]
  \item If $\sum_i\lambda_i = 1$, the point $o + \sum_i\lambda_i(a_i - o)$ does
        not depend on $o$. It is written $\sum_i\lambda_ia_i$ and called an
        \emph{affine combination}, or \emph{barycentre}, of the $a_i$.
  \item If $\sum_i\lambda_i = 0$, the vector $\sum_i\lambda_i(a_i - o)$ does
        not depend on $o$. It is also written $\sum_i\lambda_ia_i$.
  \item In all other cases, both depend on $o$.
\end{enumerate}
\end{proposition}

\begin{proof}
Replace $o$ by $o' = o + w$. By \cref{affine-spaces:prop:origin-change},
$a_i - o' = (a_i - o) - w$, so
\[
  \sum_i\lambda_i(a_i - o') = \sum_i\lambda_i(a_i - o)
  - \Bigl(\sum_i\lambda_i\Bigr)w
\]
and, since $o' = o + w$,
\[
  o' + \sum_i\lambda_i(a_i - o') = o + \sum_i\lambda_i(a_i - o)
  + \Bigl(1 - \sum_i\lambda_i\Bigr)w .
\]
The vector is the same for every $w$ if and only if $\sum_i\lambda_i = 0$,
and the point is the same for every $w$ if and only if
$\sum_i\lambda_i = 1$.
\end{proof}

\begin{intuition}
The sum of the weights is the ``number of points'' that the expression
contains. One point in total: the result is a point. Zero points in total:
the positions cancel and the result is a displacement. Any other total:
the expression still contains a multiple of the origin, and it means
nothing until the origin is fixed.
\end{intuition}

In coordinates, both notations are computed in the obvious way: the
coordinates of $\sum_i\lambda_ia_i$ are $\sum_i\lambda_ix^{(i)}$, where
$x^{(i)}$ are the coordinates of $a_i$. The proposition says when this
array of numbers is the coordinate array of a well-defined object.

\begin{example}[Meaningful and meaningless expressions]
\label{affine-spaces:ex:midpoint}
The midpoint $\frac12a + \frac12b$ and the mean $\frac1k\sum_ia_i$ are
points. The difference $b - a$ (weights $1$ and $-1$) is a vector, and so
is the second difference $a - 2b + c$. The expression $2a - b$ (weights $2$
and $-1$) is a point: the reflection of $b$ through $a$. The expressions
$a + b$ and $2a$ have total weight $2$ and are not defined. In the plane,
take a frame with origin $o$ in which $a$ and $b$ have coordinates $(1, 2)$
and $(3, 1)$, and a second origin $o'$ with coordinates $(1, 1)$. Computed
from $o$, the ``sum'' is the point of coordinates $(4, 3)$. Computed from
$o'$, the new coordinates of $a$ and $b$ are $(0, 1)$ and $(2, 0)$, their
sum is $(2, 1)$, which is the point of old coordinates $(3, 2)$. The two
results differ by $(1 - 2)w = -w$, as the proof predicts. The midpoint is
the point of old coordinates $(2, 1.5)$ in both computations
(\cref{affine-spaces:fig:sum-midpoint}).
\end{example}

\begin{figure}[ht]
\centering
\begin{tikzpicture}[>=Stealth, scale=1.25]
  % parallelogram built from the origin o
  \draw[dashed, bookaccent] (0,0) -- (1,2) -- (4,3) -- (3,1) -- cycle;
  % parallelogram built from the origin o'
  \draw[dashed, bookpitfall] (1,1) -- (1,2) -- (3,2) -- (3,1) -- cycle;
  \draw[thin, gray] (1,2) -- (3,1);
  \fill[bookaccent] (0,0) circle (1.6pt) node[below left] {$o$};
  \fill[bookpitfall] (1,1) circle (1.6pt) node[below left] {$o'$};
  \fill (1,2) circle (1.6pt) node[above left] {$a$};
  \fill (3,1) circle (1.6pt) node[below right] {$b$};
  \fill[bookaccent] (4,3) circle (1.6pt) node[above right] {$s$};
  \fill[bookpitfall] (3,2) circle (1.6pt) node[above right] {$s'$};
  \fill[bookdef] (2,1.5) circle (2pt) node[above right] {$m$};
\end{tikzpicture}
\caption{The ``sum'' of two points depends on the origin used to compute it:
$s$ is ``$a + b$'' computed from $o$ (blue parallelogram) and $s'$ is
``$a + b$'' computed from $o'$ (red parallelogram). The midpoint
$m = \frac12a + \frac12b$ (green) does not depend on the origin: it is the
centre of both parallelograms.}
\label{affine-spaces:fig:sum-midpoint}
\end{figure}

Affine combinations can be nested without leaving the family.

\begin{proposition}[Associativity of barycentres]
\label{affine-spaces:prop:associativity}
For $j = 1,\dots,m$ let $b_j = \sum_i\lambda_{ji}a_i$ be affine combinations
of the points $a_1,\dots,a_k$, and let $\mu_1,\dots,\mu_m$ be weights with
$\sum_j\mu_j = 1$. Then
\[
  \sum_j\mu_jb_j = \sum_i\Bigl(\sum_j\mu_j\lambda_{ji}\Bigr)a_i ,
\]
and the weights on the right-hand side sum to $1$.
\end{proposition}

\begin{proof}
Fix an origin $o$. By definition, $b_j - o = \sum_i\lambda_{ji}(a_i - o)$,
so
\[
  \sum_j\mu_jb_j = o + \sum_j\mu_j(b_j - o)
  = o + \sum_i\Bigl(\sum_j\mu_j\lambda_{ji}\Bigr)(a_i - o) .
\]
The total weight is
$\sum_i\sum_j\mu_j\lambda_{ji} = \sum_j\mu_j\sum_i\lambda_{ji}
= \sum_j\mu_j = 1$,
so the last expression is the affine combination claimed.
\end{proof}

\begin{definition}[Segments, convex sets, convex functions]
\label{affine-spaces:def:convex}
An affine combination with non-negative weights is a \emph{convex
combination}. The \emph{segment} $[a, b]$ is the set of points
$(1 - t)a + tb$, $t\in[0, 1]$. A subset $C\subset\Aff$ is \emph{convex} if
it contains the segment joining any two of its points. A function
$f:C\to\R$ on a convex set is \emph{convex} if
\[
  f\bigl((1 - t)a + tb\bigr)\leq(1 - t)f(a) + tf(b)
  \qquad\text{for all }a, b\in C,\ t\in[0, 1] .
\]
\end{definition}

Convexity is therefore an affine notion. It needs neither an origin nor a
way to measure lengths and angles, and it will be available long before
inner products enter in \cref{ch:inner-products}.

\begin{mlthread}[Machine learning: averages of models]
Several common techniques average parameter vectors, and all of them are
affine combinations.
\begin{itemize}[nosep]
  \item \emph{Iterate averaging}, $\bar\theta_K = \frac1K\sum_{k=1}^K\theta_k$
        (Polyak and Juditsky, 1992), and its variants along a training run
        such as stochastic weight averaging (Izmailov et al., 2018).
  \item The \emph{exponential moving average}
        $\bar\theta\leftarrow\beta\bar\theta + (1 - \beta)\theta$, with
        weights $\beta$ and $1 - \beta$.
  \item \emph{Model soups}, which average the weights of several fine-tuned
        models (Wortsman et al., 2022), and linear interpolation
        $(1 - t)\theta_0 + t\theta_1$ between two models.
  \item \emph{Task arithmetic} (Ilharco et al., 2023) forms
        $\theta_0 + \sum_i\lambda_i(\theta_i - \theta_0)$, where $\theta_0$
        are pretrained weights and $\theta_i$ fine-tuned ones. The
        differences $\theta_i - \theta_0$ are called task vectors, and
        rightly so: they are vectors, and a point plus a combination of
        vectors is a point whatever the $\lambda_i$. The sum of the weights
        on the points is $1 - \sum_i\lambda_i + \sum_i\lambda_i = 1$.
\end{itemize}
On the data side, \emph{mixup} (Zhang et al., 2018) trains on convex
combinations $(1 - t)x_i + tx_j$ of pairs of inputs, and centring a data
set replaces the points $x_i$ by the vectors $x_i - \bar x$, where
$\bar x = \frac1n\sum_ix_i$ is their mean. A covariance matrix is built
from these vectors, which is why it does not depend on the origin.
\end{mlthread}

\section{Affine subspaces}\label{affine-spaces:sec:subspaces}

\begin{definition}[Affine subspace]\label{affine-spaces:def:subspace}
A non-empty subset $\Bff\subset\Aff$ is an \emph{affine subspace} if there
are a point $b\in\Bff$ and a linear subspace $W\subset V$ such that
\[
  \Bff = b + W = \{b + w : w\in W\} .
\]
\end{definition}

\begin{proposition}\label{affine-spaces:prop:subspace}
Let $\Bff = b + W$ be an affine subspace of $\Aff$.
\begin{enumerate}[nosep]
  \item For every $b'\in\Bff$, $\Bff = b' + W$, and
        $W = \{c - b' : c\in\Bff\}$. In particular $W$ is determined by
        $\Bff$; it is called the \emph{direction} of $\Bff$.
  \item With the operation of $\Aff$ restricted to $\Bff\times W$, $\Bff$ is
        an affine space with direction $W$.
\end{enumerate}
\end{proposition}

\begin{proof}
(1) Write $b' = b + w_0$ with $w_0\in W$. For $w\in W$,
$b' + w = b + (w_0 + w)\in b + W$ and $b + w = b' + (w - w_0)\in b' + W$,
because $W$ is a linear subspace. So $b + W = b' + W$. Then
$\{c - b' : c\in\Bff\} = \{(b' + w) - b' : w\in W\} = W$.

(2) By (1), $b' + w\in\Bff$ for all $b'\in\Bff$ and $w\in W$, so the
operation is well defined. Axioms 1 and 2 hold in $\Bff$ because they hold
in $\Aff$. For axiom 3, let $b', c\in\Bff$. By (1), $c - b'\in W$, and it is
the only vector of $V$, hence of $W$, that leads from $b'$ to $c$.
\end{proof}

Affine subspaces of dimension $0$, $1$ and $n - 1$ are the points, the
\emph{lines} and the \emph{hyperplanes} of $\Aff$. Two affine subspaces with
the same direction are \emph{parallel}
(\cref{affine-spaces:fig:subspace}). The spaces $S$ and $H$ of
\cref{affine-spaces:ex:system,affine-spaces:ex:simplex} are affine
subspaces of $\R^n$.

\begin{figure}[ht]
\centering
\begin{tikzpicture}[>=Stealth, scale=1.1]
  \draw[thick, bookaccent] (-0.6,-0.3) -- (5.4,2.7)
    node[right] {$\Bff = b + W$};
  \draw[thick, bookaccent!50] (0.4,-1.5) -- (6.4,1.5)
    node[right] {$\Bff' = b' + W$};
  \fill (1,0.5) circle (1.6pt) node[above left] {$b$};
  \fill (3,1.5) circle (1.6pt) node[above left] {$c$};
  \draw[->, very thick, bookdef] (1,0.5) -- (3,1.5)
    node[midway, below right] {$c - b\in W$};
  \fill (2,-0.7) circle (1.6pt) node[below right] {$b'$};
\end{tikzpicture}
\caption{An affine subspace is a point plus a linear subspace of
displacements. Two parallel lines share the direction $W$; neither of them
is a vector space, and neither contains a preferred point.}
\label{affine-spaces:fig:subspace}
\end{figure}

\begin{proposition}[Characterization by affine combinations]
\label{affine-spaces:prop:subspace-comb}
A non-empty subset $\Bff\subset\Aff$ is an affine subspace if and only if
every affine combination of points of $\Bff$ belongs to $\Bff$.
\end{proposition}

\begin{proof}
Suppose $\Bff = b + W$ and let $b_i = b + w_i$ be points of $\Bff$ and
$\sum_i\lambda_i = 1$. Computing with the origin $b$,
$\sum_i\lambda_ib_i = b + \sum_i\lambda_iw_i$, which is in $b + W$.

Conversely, suppose $\Bff$ is stable under affine combinations, fix
$b\in\Bff$ and let $W = \{c - b : c\in\Bff\}$, so that $\Bff = b + W$. We
show that $W$ is a linear subspace. It contains $0 = b - b$. If
$w = c - b\in W$ and $t\in\R$, the point $(1 - t)b + tc = b + tw$ is in
$\Bff$, so $tw\in W$. If $w_1 = c_1 - b$ and $w_2 = c_2 - b$ are in $W$, the
point $c_1 + c_2 - b$ (weights $1$, $1$, $-1$) is in $\Bff$, and computed
with the origin $b$ it equals $b + w_1 + w_2$, so $w_1 + w_2\in W$.
\end{proof}

\begin{definition}[Affine hull, convex hull]
The \emph{affine hull} $\aff(S)$ of a non-empty set $S\subset\Aff$ is the
set of all affine combinations of points of $S$, and its \emph{convex hull}
$\conv(S)$ is the set of all convex combinations of points of $S$.
\end{definition}

By \cref{affine-spaces:prop:associativity}, an affine combination of
affine combinations of points of $S$ is again one, so $\aff(S)$ is an
affine subspace by \cref{affine-spaces:prop:subspace-comb}. It is the
smallest one containing $S$, since any affine subspace containing $S$
contains all affine combinations of its points. Likewise $\conv(S)$ is
convex, because the weights produced by
\cref{affine-spaces:prop:associativity} are non-negative when the
$\lambda_{ji}$ and the $\mu_j$ are. It is the smallest convex set containing
$S$: a convex set $C$ contains every convex combination
$\sum_{i=1}^k\lambda_ia_i$ of its points, by induction on $k$. For $k = 2$
this is the definition. For $k > 2$ with $\lambda_k\neq1$, the point
$b = \sum_{i<k}\frac{\lambda_i}{1 - \lambda_k}a_i$ is a convex combination of
$k - 1$ points, hence in $C$, and
$\sum_{i\leq k}\lambda_ia_i = (1 - \lambda_k)b + \lambda_ka_k$ by
\cref{affine-spaces:prop:associativity}, which is on the segment
$[b, a_k]$. For instance, the hyperplane $H$ of
\cref{affine-spaces:ex:simplex} is the affine hull of the points
$e_1,\dots,e_n$ of $\R^n$, and the set of probability vectors is their
convex hull, the \emph{simplex}.

\section{Affine maps}\label{affine-spaces:sec:maps}

Let $\Aff$ and $\Bff$ be affine spaces with directions $V$ and $W$.

\begin{definition}[Affine map]\label{affine-spaces:def:map}
A map $f:\Aff\to\Bff$ is \emph{affine} if there is a linear map
$\lin f:V\to W$ such that
\[
  f(a + v) = f(a) + \lin f(v)\qquad\text{for all }a\in\Aff,\ v\in V .
\]
The map $\lin f$ is the \emph{linear part} of $f$.
\end{definition}

The linear part is determined by $f$, since
$\lin f(v) = f(a + v) - f(a)$. An affine map sends points to points, and
its linear part says what it does to displacements: if $b = a + v$ then
$f(b) - f(a) = \lin f(b - a)$.

\begin{proposition}[One base point is enough]
\label{affine-spaces:prop:one-point}
Let $f:\Aff\to\Bff$ be a map. If there are a point $o\in\Aff$ and a linear
map $L:V\to W$ such that $f(o + v) = f(o) + L(v)$ for all $v\in V$, then
$f$ is affine with linear part $L$.
\end{proposition}

\begin{proof}
Let $a = o + u$ be any point. By axiom 2 and the linearity of $L$,
\[
  f(a + v) = f\bigl(o + (u + v)\bigr) = f(o) + L(u) + L(v)
  = f(o + u) + L(v) = f(a) + L(v) .
\qedhere
\]
\end{proof}

\paragraph{In coordinates.} Choose frames $(o; e_1,\dots,e_n)$ of $\Aff$ and
$(o'; e'_1,\dots,e'_m)$ of $\Bff$. Let $A$ be the matrix of $\lin f$ in the
two bases and let $\beta$ be the coordinates of the point $f(o)$. If $a$
has coordinates $x$, then $f(a) = f(o) + \lin f(a - o)$ has coordinates
\[
  y = Ax + \beta .
\]
Conversely, by \cref{affine-spaces:prop:one-point}, every map of this form
is affine. Now move the origin of $\Aff$ by a vector of components
$\omega$: the new coordinates are $x' = x - \omega$, and
\[
  y = A(x' + \omega) + \beta = Ax' + (\beta + A\omega) .
\]
The matrix has not changed; the constant term has.

\begin{keyidea}
In $y = Ax + \beta$, the matrix $A$ represents the linear part of the map
and depends only on the two bases. The constant $\beta$ is the position of
the image of the origin and depends on both origins. An affine map has no
intrinsic ``constant term''.
\end{keyidea}

\begin{proposition}[Affine maps and affine combinations]
\label{affine-spaces:prop:map-comb}
A map $f:\Aff\to\Bff$ is affine if and only if it preserves affine
combinations: $f\bigl(\sum_i\lambda_ia_i\bigr) = \sum_i\lambda_if(a_i)$
whenever $\sum_i\lambda_i = 1$.
\end{proposition}

\begin{proof}
Suppose $f$ is affine and $\sum_i\lambda_i = 1$. Fix $o\in\Aff$. Then
\[
  f\Bigl(o + \sum_i\lambda_i(a_i - o)\Bigr)
  = f(o) + \sum_i\lambda_i\lin f(a_i - o)
  = f(o) + \sum_i\lambda_i\bigl(f(a_i) - f(o)\bigr) ,
\]
which is the affine combination $\sum_i\lambda_if(a_i)$ computed in $\Bff$
with the origin $f(o)$.

Conversely, suppose $f$ preserves affine combinations. Fix $o\in\Aff$ and
define $L(v) = f(o + v) - f(o)$. For $t\in\R$, the point $o + tv$ is the
affine combination $(1 - t)o + t(o + v)$, so
$f(o + tv) = (1 - t)f(o) + tf(o + v) = f(o) + tL(v)$, hence
$L(tv) = tL(v)$. For $u, v\in V$, the point $o + u + v$ is the affine
combination $(o + u) + (o + v) - o$, so
$f(o + u + v) = f(o + u) + f(o + v) - f(o) = f(o) + L(u) + L(v)$, hence
$L(u + v) = L(u) + L(v)$. So $L$ is linear, and $f$ is affine by
\cref{affine-spaces:prop:one-point}.
\end{proof}

\begin{proposition}[Composition, inverse, translations]
\label{affine-spaces:prop:composition}
\begin{enumerate}[nosep]
  \item If $f:\Aff\to\Bff$ and $g:\Bff\to\mathcal C$ are affine, then
        $g\circ f$ is affine, with linear part $\lin g\circ\lin f$.
  \item An affine map $f$ is bijective if and only if $\lin f$ is, and then
        $f^{-1}$ is affine with linear part $(\lin f)^{-1}$.
  \item For $u\in V$, the \emph{translation} $\tau_u:a\mapsto a + u$ is
        affine with linear part $\id_V$. Conversely, an affine map
        $f:\Aff\to\Aff$ with linear part $\id_V$ is a translation.
\end{enumerate}
\end{proposition}

\begin{proof}
(1) $g(f(a + v)) = g\bigl(f(a) + \lin f(v)\bigr)
= g(f(a)) + \lin g(\lin f(v))$.

(2) Fix $o\in\Aff$. Since $\Theta_o$ and $\Theta_{f(o)}$ are bijections and
$\Theta_{f(o)}\circ f = \lin f\circ\Theta_o$, the map $f$ is bijective if
and only if $\lin f$ is. In that case, for $b\in\Bff$ and $w\in W$, apply
$f$ to the point $f^{-1}(b) + (\lin f)^{-1}(w)$: the result is $b + w$.
Hence $f^{-1}(b + w) = f^{-1}(b) + (\lin f)^{-1}(w)$.

(3) By axiom 2, $\tau_u(a + v) = a + v + u = \tau_u(a) + v$. Conversely, if
$\lin f = \id_V$ then, by rule (4) of \cref{affine-spaces:prop:rules},
$f(a + v) - (a + v) = (f(a) + v) - (a + v) = f(a) - a$ for all $v$. So the
vector $u = f(a) - a$ is the same for all points, and $f = \tau_u$.
\end{proof}

\begin{proposition}[Level sets of affine maps]
\label{affine-spaces:prop:fibre}
Let $f:\Aff\to\Bff$ be affine and $c\in\Bff$. The set
$f^{-1}(c) = \{a\in\Aff : f(a) = c\}$ is either empty or an affine subspace
of $\Aff$ with direction $\ker\lin f$.
\end{proposition}

\begin{proof}
Suppose $f(a_0) = c$. For $v\in V$, $f(a_0 + v) = c + \lin f(v)$, and by
rule (1) of \cref{affine-spaces:prop:rules} this equals $c$ if and only if
$\lin f(v) = 0$. So $f^{-1}(c) = a_0 + \ker\lin f$.
\end{proof}

This is the familiar statement that the general solution of $Ax = b$ is a
particular solution plus the general solution of $Ax = 0$
(\cref{affine-spaces:ex:system}).

\begin{proposition}[Fixed points]\label{affine-spaces:prop:fixed-point}
Let $f:\Aff\to\Aff$ be affine, with $\Aff$ finite-dimensional. If $1$ is
not an eigenvalue of $\lin f$, then $f$ has exactly one fixed point $c$,
and
\[
  f(c + v) = c + \lin f(v)\quad\text{for all }v\in V,
  \qquad\text{so that}\qquad
  f^k(a) - c = (\lin f)^k(a - c)
\]
for every point $a$ and every $k\geq0$.
\end{proposition}

\begin{proof}
Fix $o\in\Aff$. The point $o + v$ is fixed if and only if
$f(o) + \lin f(v) = o + v$, that is, $(\id_V - \lin f)(v) = f(o) - o$. Since
$1$ is not an eigenvalue of $\lin f$, the map $\id_V - \lin f$ is injective,
hence bijective in finite dimension, and this equation has exactly one
solution. If $c$ is the fixed point, $f(c + v) = f(c) + \lin f(v)
= c + \lin f(v)$, and the formula for $f^k$ follows by induction.
\end{proof}

\begin{intuition}
Seen from its fixed point, an affine map is linear. The position of a
point relative to $c$ is a vector, and the map transforms this vector by
$\lin f$. All questions about the iterates of $f$ become questions about
the powers of a linear map.
\end{intuition}

\begin{example}[Gradient descent on a quadratic is an affine map]
\label{affine-spaces:ex:quadratic}
In coordinates, let $L(x) = \frac12x^\top Hx - b^\top x$ with $H$ symmetric
and invertible. The vector of its partial derivatives is $Hx - b$, and one
step of gradient descent with step size $\eta$ is
\[
  x\ \longmapsto\ x - \eta(Hx - b) = (I - \eta H)x + \eta b ,
\]
an affine map with linear part $I - \eta H$. Its fixed point is the
solution $x^*$ of $Hx = b$, and by \cref{affine-spaces:prop:fixed-point},
\[
  x_k - x^* = (I - \eta H)^k(x_0 - x^*) .
\]
The \emph{error} $x_k - x^*$ is a vector and evolves linearly. Whether and
how fast gradient descent converges on this loss is entirely a question
about the powers of $I - \eta H$, which \cref{ch:spectral-svd} answers.
\end{example}

\begin{mlthread}[Machine learning: layers are affine, not linear]
A ``linear layer'' $x\mapsto Wx + b$ is an affine map. Three consequences
follow from this section.
\begin{itemize}[nosep]
  \item The composition of affine maps is affine
        (\cref{affine-spaces:prop:composition}), with
        $W_2(W_1x + b_1) + b_2 = (W_2W_1)x + (W_2b_1 + b_2)$. A stack of such
        layers without non-linearities computes nothing more than a single
        one.
  \item Translating the inputs, for instance subtracting their mean,
        changes the bias and not the weight matrix: $Wx + b
        = W(x - \mu) + (b + W\mu)$. A layer with a bias can absorb any
        choice of origin of its input space. A layer without bias cannot:
        the family of maps $x\mapsto Wx$ singles out the input that is
        sent to zero (\cref{affine-spaces:exr:6}).
  \item If $g$ is affine and $f$ is convex, then $f\circ g$ is convex
        (\cref{affine-spaces:exr:9}). This is why least squares and
        logistic regression are convex problems in their parameters: the
        predictions depend on the parameters through an affine map.
\end{itemize}
\end{mlthread}

\section{Points at height one}\label{affine-spaces:sec:homogeneous}

There is a way of storing points and vectors in arrays that makes the type
checking automatic. Append one extra number to each array:
\[
  \text{a point with coordinates }x\in\R^n\ \longmapsto\
  \begin{bmatrix}x\\1\end{bmatrix},
  \qquad
  \text{a vector with coordinates }v\in\R^n\ \longmapsto\
  \begin{bmatrix}v\\0\end{bmatrix} .
\]

\begin{proposition}[Homogeneous coordinates]
\label{affine-spaces:prop:homogeneous}
With this convention:
\begin{enumerate}[nosep]
  \item the last coordinate of $\sum_i\lambda_i\begin{bsmallmatrix}x^{(i)}\\
        1\end{bsmallmatrix}$ is $\sum_i\lambda_i$, so a weighted sum of
        points is a point when the last coordinate is $1$, a vector when it
        is $0$, and neither otherwise;
  \item the affine map $y = Ax + \beta$ acts on points and its linear part
        acts on vectors through the same $(m + 1)\times(n + 1)$ matrix,
        \[
          \begin{bmatrix}A & \beta\\0 & 1\end{bmatrix}
          \begin{bmatrix}x\\1\end{bmatrix}
          = \begin{bmatrix}Ax + \beta\\1\end{bmatrix},
          \qquad
          \begin{bmatrix}A & \beta\\0 & 1\end{bmatrix}
          \begin{bmatrix}v\\0\end{bmatrix}
          = \begin{bmatrix}Av\\0\end{bmatrix} ;
        \]
  \item composition of affine maps corresponds to the product of these
        matrices.
\end{enumerate}
\end{proposition}

\begin{proof}
(1) and (2) are direct computations. For (3),
\[
  \begin{bmatrix}A_2 & \beta_2\\0 & 1\end{bmatrix}
  \begin{bmatrix}A_1 & \beta_1\\0 & 1\end{bmatrix}
  = \begin{bmatrix}A_2A_1 & A_2\beta_1 + \beta_2\\0 & 1\end{bmatrix},
\]
which is the matrix of $x\mapsto A_2(A_1x + \beta_1) + \beta_2$.
\end{proof}

Geometrically, the $n$-dimensional affine space has been placed in
$\R^{n+1}$ as the hyperplane at height $1$, and its direction as the
parallel hyperplane through the origin, at height $0$
(\cref{affine-spaces:fig:homogeneous}). The height of an array is the
total weight of \cref{affine-spaces:prop:affine-comb}, and the constant
$\beta$ is invisible to vectors because it is multiplied by their last
coordinate, $0$. In machine learning this device is known as the
\emph{bias trick}: appending a constant feature $1$ to the input turns the
affine layer $Wx + b$ into the linear map $[W\ b]$ of the extended input.

\begin{figure}[ht]
\centering
\begin{tikzpicture}[>=Stealth, scale=1.0]
  \draw[->, gray] (-3.4,0) -- (5.6,0);
  \draw[->, gray] (0,-0.4) -- (0,2.7) node[above] {height};
  \draw[thick, bookaccent] (-3.4,1) -- (5.6,1)
    node[right] {points};
  \draw[thick, bookdef] (-3.4,0) -- (5.6,0) node[right] {vectors};
  \draw[thin, gray, dashed] (-3.4,2) -- (5.6,2);
  \node[gray, anchor=west] at (5.6,2) {height $2$};
  \fill (1,1) circle (1.6pt) node[below] {$a$};
  \fill (3,1) circle (1.6pt) node[below] {$b$};
  \fill[bookpitfall] (4,2) circle (1.6pt) node[above] {$a + b$};
  \fill[bookaccent] (2,1) circle (1.8pt) node[above] {$\frac12a + \frac12b$};
  \draw[->, very thick, bookdef] (0,0) -- (2,0)
    node[midway, below] {$b - a$};
  \fill[bookaccent] (-1,1) circle (1.6pt) node[above] {$2a - b$};
\end{tikzpicture}
\caption{Homogeneous coordinates for a one-dimensional affine space. Points
are stored at height $1$ and vectors at height $0$. The midpoint and the
reflection $2a - b$ stay at height $1$; the difference $b - a$ falls to
height $0$; the ``sum'' $a + b$ lands at height $2$, where nothing lives.}
\label{affine-spaces:fig:homogeneous}
\end{figure}

\section{The origin test}\label{affine-spaces:sec:origin-test}

The computer does not check heights. On arrays, \texttt{x + y} and
\texttt{2 * x} are as easy to write as \texttt{x + v}, and a formula that
uses them still returns numbers. The question is whether those numbers
mean anything.

\begin{definition}[Origin test]\label{affine-spaces:def:origin-test}
A formula computed on coordinates \emph{passes the origin test} if its
result does not depend on the origin of the frame: when the origin is moved
by $w$, the result has its coordinates shifted by $-\omega$ if it is meant
to be a point, and unchanged if it is meant to be a vector or a number.
\end{definition}

By \cref{affine-spaces:prop:origin-change}, moving the origin by $w$ means
replacing every array $x$ that stores a point by $x - \omega$ and leaving
every array that stores a vector unchanged. \Cref{affine-spaces:tab:test}
applies the test to common expressions. The last two lines anticipate
\cref{ch:inner-products}: the distance between two points is the length of
a vector, whereas the ``norm of a point'' and the cosine similarity
$x\cdot y/(\norm x\norm y)$ of two points are computed on $x - o$ and
$y - o$, and change with $o$. For instance, the points of coordinates
$(1, 0)$ and $(0, 1)$ have cosine similarity $0$; seen from the origin
placed at $(-1, -1)$ their coordinates are $(2, 1)$ and $(1, 2)$, and the
cosine similarity is $\frac45$. Their distance is $\sqrt2$ in both cases.

\begin{table}[ht]
\centering
\begin{tabular}{@{}lll@{}}
\toprule
Expression on coordinates & Meant to be & Origin test \\
\midrule
$x + v$ & a point & passes \\
$y - x$ & a vector & passes \\
$\frac12(x + y)$,\quad $\frac1k\sum_ix^{(i)}$ & a point & passes \\
$\beta\bar x + (1 - \beta)x$ & a point & passes \\
$x + \beta(x - x_{\mathrm{prev}})$ & a point & passes \\
$x + y$,\quad $2x$,\quad $-x$ & a point & fails \\
$(1 - \eta\lambda)x$ & a point & fails \\
$\norm{y - x}$ & a number & passes \\
$\norm x$,\quad $x\cdot y$,\quad cosine similarity of $x$ and $y$
  & a number & fails \\
\bottomrule
\end{tabular}
\caption{The origin test on common expressions; $x$, $y$, $\bar x$ store
points and $v$ stores a vector.}
\label{affine-spaces:tab:test}
\end{table}

\paragraph{Update rules.} A training algorithm produces at each step a
vector to add to the current point. Abstractly, it uses a map
$g:\Aff\to V$, a \emph{vector field}, and iterates
\begin{equation}\label{affine-spaces:eq:update}
  a_{k+1} = a_k + \eta\,g(a_k) .
\end{equation}
For gradient descent, $g = -\nabla L$; how a loss function $L$ produces
such a field, and what else that requires, is the subject of
\cref{ch:differential,ch:gradient-hessian}. Whatever $g$ is,
\eqref{affine-spaces:eq:update} is a point plus a vector, and it passes the
test: if $x_k$ are the coordinates of $a_k$ and $\gamma(x)$ those of
$g(a)$, then in the frame with origin moved by $w$ the same field reads
$\gamma'(x') = \gamma(x' + \omega)$, and
$x'_{k+1} = x'_k + \eta\gamma'(x'_k)$ is $x_{k+1} - \omega$ whenever
$x'_k = x_k - \omega$. The same holds for the heavy-ball momentum update,
which adds the vector $\beta(a_k - a_{k-1})$, and for the averages of
\cref{affine-spaces:sec:combinations}.

\begin{example}[Weight decay has a hidden origin]
\label{affine-spaces:ex:decay}
Weight decay modifies the update \eqref{affine-spaces:eq:update} into
\[
  x_{k+1} = (1 - \eta\lambda)x_k + \eta\gamma(x_k) .
\]
The coefficient of the point $x_k$ is not $1$, and the formula fails the
test: computed in two frames whose origins differ by $w$, from the same
point, it gives two points that differ by the vector $\eta\lambda w$
(\cref{affine-spaces:exr:5}). Written without coordinates, with $o$ the
point whose coordinates are zero, the update is
\[
  a_{k+1} = \bigl(o + (1 - \eta\lambda)(a_k - o)\bigr) + \eta\,g(a_k) :
\]
a contraction of ratio $1 - \eta\lambda$ towards $o$, followed by the usual
step. Each frame contracts towards its own origin.

To see the effect on the result, take $g(a) = c - a$, a field that pulls
every point towards a point $c$ (it is $-\nabla L$ for
$L(a) = \frac12\norm{a - c}^2$ with the dot product). Then
$a_{k+1} = a_k + \eta\bigl((c - a_k) + \lambda(o - a_k)\bigr)$ is an affine
map with linear part $(1 - \eta(1 + \lambda))\id_V$. For
$0 < \eta < 2/(1 + \lambda)$ it converges, by
\cref{affine-spaces:prop:fixed-point}, to its fixed point
\[
  a^* = \frac1{1 + \lambda}\,c + \frac\lambda{1 + \lambda}\,o ,
\]
a barycentre of $c$ and $o$. With $\lambda = 1$ and $c$ of coordinates
$(3, 1)$: if $o$ is the point $(0, 0)$, training ends at $(1.5, 0.5)$; if
the decay is centred at the point $(2, 2)$ instead, it ends at
$(2.5, 1.5)$ (\cref{affine-spaces:fig:decay}).
\end{example}

\begin{figure}[ht]
\centering
\begin{tikzpicture}[>=Stealth, scale=1.5]
  \draw[thin, gray!60] (0,0) -- (3,1) -- (2,2);
  \fill (3,1) circle (1.6pt) node[right] {$c$};
  \fill[bookaccent] (0,0) circle (1.6pt) node[below left] {$o$};
  \fill[bookpitfall] (2,2) circle (1.6pt) node[above right] {$o'$};
  \fill (0.5,2.5) circle (1.6pt) node[above left] {$a_0$};
  \draw[bookaccent, thick, mark=*, mark size=1.2pt] plot coordinates
    {(0.5,2.5) (1.0,1.5) (1.25,1.0) (1.375,0.75) (1.4375,0.625)
     (1.46875,0.5625) (1.5,0.5)};
  \draw[bookpitfall, thick, mark=*, mark size=1.2pt] plot coordinates
    {(0.5,2.5) (1.5,2.0) (2.0,1.75) (2.25,1.625) (2.375,1.5625)
     (2.4375,1.53125) (2.5,1.5)};
  \node[bookaccent, anchor=north west] at (1.5,0.45) {$a^*$ (decay to $o$)};
  \node[bookpitfall, anchor=west] at (2.56,1.55) {$a^*$ (decay to $o'$)};
\end{tikzpicture}
\caption{The same training run with weight decay centred at two different
points ($\lambda = 1$, $\eta = 0.25$, field $g(a) = c - a$). Each run ends
at the midpoint of $c$ and its own centre of decay.}
\label{affine-spaces:fig:decay}
\end{figure}

\begin{remark}[Coupled and decoupled decay: AdamW]
\label{affine-spaces:rem:adamw}
With plain gradient descent, adding $-\lambda x$ to the field and
multiplying $x$ by $1 - \eta\lambda$ are the same thing:
$x + \eta\bigl(\gamma(x) - \lambda x\bigr)
= (1 - \eta\lambda)x + \eta\gamma(x)$. With an adaptive optimizer such as
Adam they are not. In the \emph{coupled} form, the term $\lambda x$ is added
to the gradient and goes through the coordinatewise rescaling of the
optimizer. In the \emph{decoupled} form, AdamW (Loshchilov and Hutter,
2019), the optimizer computes its step $s_k$ from the gradient alone, and
the update is
\[
  x_{k+1} = (1 - \eta\lambda)x_k + s_k .
\]
Decoupling changes how the decay interacts with the rescaling. It does not
change the geometry. The step $s_k$ of Adam is a vector that passes the
origin test, and the factor $1 - \eta\lambda$ is the same contraction
towards the point whose coordinates are zero:
$a_{k+1} = \bigl(o + (1 - \eta\lambda)(a_k - o)\bigr) + s_k$. Both forms
fail the origin test, and the decoupled form shows why in its formula. To
centre the decay of AdamW at pretrained weights $\theta_0$, one stores
$\theta - \theta_0$ as the parameter, or applies the contraction by hand.
\end{remark}

\begin{pitfall}[What the test says, and what it does not]
The origin test does not say that weight decay is wrong. It says that the
formula contains a choice, the point towards which the parameters are
pulled, and that the notation $\lambda\theta$ hides it. For a network
trained from scratch the zero weights are a meaningful centre: the
parameter space of a neural network is a genuine vector space, in which
the zero weights are the network whose output ignores its input. For a
network fine-tuned from pretrained weights $\theta_0$, pulling towards
$\theta_0$, that is, replacing $\lambda\theta$ by $\lambda(\theta - \theta_0)$,
is often the better choice (the L2-SP regularizer of Li, Grandvalet and
Davoine, 2018). The test sorts formulas into those that use an origin and
those that do not, so that the choice is made knowingly.
\end{pitfall}

Moving the origin is one of two natural changes of frame. The other one,
changing the basis, is the subject of the \emph{basis test} of
\cref{ch:duality}. Coordinatewise optimizers such as Adam pass the origin
test and fail the basis test.

\section{Where the chapter leads}\label{affine-spaces:sec:outlook}

An affine space has points, displacements, lines, parallelism, barycentres
and convexity. It has no lengths and no angles, so no notion of
``steepest''. The following chapters add structure one layer at a time.
\begin{itemize}[nosep]
  \item \Cref{ch:duality} studies the vector space $V$ and its dual $V^*$,
        the space of linear forms on $V$. An \emph{affine form}
        $\ell:\Aff\to\R$ has a linear part $\lin\ell\in V^*$: the first
        example of a covector, and the model for the differential of a
        loss.
  \item \Cref{ch:inner-products} adds an inner product on $V$, which makes
        $\Aff$ a \emph{Euclidean affine space} and turns covectors into
        vectors.
  \item \Cref{ch:differential} defines the differential of a map
        $f:\Aff\to\Bff$ at a point $a$ as the linear part of the affine map
        that best approximates $f$ near $a$. The velocity of a curve
        $\gamma$ in $\Aff$,
        \[
          \gamma'(t) = \lim_{h\to0}\frac{\gamma(t + h) - \gamma(t)}h ,
        \]
        is a limit of vectors of $V$, whatever the point $\gamma(t)$. In an
        affine space, velocities at all points live in the same space $V$.
        This is what allows a vector computed at $a_k$ to be added to
        $a_k$, and it is exactly what is lost on a curved space
        (\cref{ch:submanifolds}).
\end{itemize}

\paragraph{Sources.} The two books that inspired this one treat the
subject differently. Hubbard and Hubbard (\emph{Vector Calculus, Linear
Algebra, and Differential Forms}, 5th ed., \S1.1) distinguish points from
vectors, with separate notations and the rules that points cannot be added
while vectors can, without naming the structure; the temperature example
is adapted from theirs. Loomis and Sternberg (\emph{Advanced Calculus},
\S1.2) build the
correspondence between Euclidean space and $\R^3$ from a chosen origin and
then transfer the geometric language to the vector space; affine subspaces
appear in their \S1.4 as translates of linear subspaces of a vector space.
Systematic treatments of affine spaces are in Berger (\emph{Geometry~I},
Springer, 1987, Chapter~2) and Gallier (\emph{Geometric Methods and
Applications}, 2nd ed., Springer, 2011, Chapter~2). The machine learning
references are: B.~Polyak and A.~Juditsky, Acceleration of stochastic
approximation by averaging, \emph{SIAM J.\ Control Optim.}, 1992;
P.~Izmailov et al., Averaging weights leads to wider optima and better
generalization, UAI, 2018; M.~Wortsman et al., Model soups, ICML, 2022;
G.~Ilharco et al., Editing models with task arithmetic, ICLR, 2023;
H.~Zhang et al., mixup: Beyond empirical risk minimization, ICLR, 2018;
X.~Li, Y.~Grandvalet and F.~Davoine, Explicit inductive bias for transfer
learning with convolutional networks, ICML, 2018; I.~Loshchilov and
F.~Hutter, Decoupled weight decay regularization, ICLR, 2019.

\paragraph{Notebook.} \notebook{01-affine-spaces.ipynb} checks the
numerical claims of this chapter with PyTorch.

\begin{chaptersummary}
\begin{itemize}[nosep]
  \item An affine space is a set of points on which a vector space of
        displacements acts: point $+$ vector $=$ point, point $-$ point $=$
        vector, with exactly one displacement between any two points.
        There is no sum of points and no origin.
  \item Choosing an origin represents points by vectors. When the origin
        moves by $w$, coordinates of points change by $-w$; coordinates of
        vectors do not change.
  \item A weighted sum of points is a point if the weights sum to $1$, a
        vector if they sum to $0$, and meaningless otherwise. Midpoints,
        means, moving averages and convexity are affine notions.
  \item Affine subspaces are translates $b + W$ of linear subspaces;
        solution sets of linear systems are the typical example.
  \item An affine map is $x\mapsto Ax + \beta$ in coordinates. The matrix
        $A$ is its linear part and acts on displacements; $\beta$ depends
        on the origins. Seen from a fixed point, an affine map is linear.
  \item Storing points with a last coordinate $1$ and vectors with a last
        coordinate $0$ makes these rules automatic (the bias trick).
  \item The origin test detects formulas that depend on the origin.
        Updates of the form point $+$ vector pass; weight decay fails,
        with SGD as with AdamW, because it pulls the parameters towards a
        chosen point.
\end{itemize}
\end{chaptersummary}

\section*{Exercises}
\addcontentsline{toc}{section}{Exercises}

\begin{exercise}\label{affine-spaces:exr:1}
Let $a, b, c, d$ be points of an affine space. Show that $b - a = c - d$
if and only if $d - a = c - b$ (the quadrilateral $abcd$ is then a
parallelogram), and that this happens if and only if the midpoints of
$[a, c]$ and $[b, d]$ coincide.
\solutionlink{affine-spaces:exr:1}
\end{exercise}

\begin{exercise}\label{affine-spaces:exr:2}
Let $a, b, c$ be points and $u$ a vector. Say whether each expression is a
point, a vector, or undefined:
\[
  a + 2u,\quad \tfrac12(a + b),\quad a + b - c,\quad a - 2b + c,\quad
  2a - b,\quad a + b,\quad \tfrac13(a + b + c) - a .
\]
\solutionlink{affine-spaces:exr:2}
\end{exercise}

\begin{exercise}\label{affine-spaces:exr:3}
Show that affine combinations of \emph{two} points are enough in
\cref{affine-spaces:prop:subspace-comb,affine-spaces:prop:map-comb}:
\begin{enumerate}[nosep]
  \item a non-empty subset $\Bff\subset\Aff$ that contains the line through
        any two of its points is an affine subspace;
  \item a map $f$ such that $f((1 - t)a + tb) = (1 - t)f(a) + tf(b)$ for
        all points $a, b$ and all $t\in\R$ is affine.
\end{enumerate}
\hint{$o + u + v$ is the midpoint of $o + 2u$ and $o + 2v$.}
\solutionlink{affine-spaces:exr:3}
\end{exercise}

\begin{exercise}\label{affine-spaces:exr:4}
Let $a, b, c$ be three points and $g = \frac13a + \frac13b + \frac13c$.
Use \cref{affine-spaces:prop:associativity} to show that $g$ lies on the
segment joining $a$ to the midpoint of $[b, c]$, two thirds of the way
from $a$. Conclude that the three medians of a triangle meet at $g$.
\solutionlink{affine-spaces:exr:4}
\end{exercise}

\begin{exercise}\label{affine-spaces:exr:5}
Weight decay is run in coordinates $x$ and in the coordinates
$x' = x - \omega$ of a frame whose origin has been moved by $w$:
$x\mapsto(1 - \eta\lambda)x + \eta\gamma(x)$ and
$x'\mapsto(1 - \eta\lambda)x' + \eta\gamma'(x')$, with
$\gamma'(x') = \gamma(x' + \omega)$.
\begin{enumerate}[nosep]
  \item Starting from the same point, by how much do the two new points
        differ after one step?
  \item For $\gamma\equiv0$ and $0 < \eta\lambda < 2$, where do the two
        sequences converge?
  \item For $\gamma(x) = \chi - x$ (the field of
        \cref{affine-spaces:ex:decay}, with $\chi$ the coordinates of $c$),
        verify the fixed point and the condition on $\eta$ given in that
        example.
\end{enumerate}
\solutionlink{affine-spaces:exr:5}
\end{exercise}

\begin{exercise}\label{affine-spaces:exr:6}
Let $\mathcal F_{\mathrm{aff}}$ be the set of maps $\R^n\to\R^m$ of the form
$x\mapsto Wx + b$ and $\mathcal F_{\mathrm{lin}}$ the set of maps
$x\mapsto Wx$. For $\omega\in\R^n$ let $\tau_\omega(x) = x + \omega$.
\begin{enumerate}[nosep]
  \item Show that $f\circ\tau_\omega\in\mathcal F_{\mathrm{aff}}$ for every
        $f\in\mathcal F_{\mathrm{aff}}$, and give its weight matrix and
        bias.
  \item Show that for $f\in\mathcal F_{\mathrm{lin}}$,
        $f\circ\tau_\omega\in\mathcal F_{\mathrm{lin}}$ if and only if
        $W\omega = 0$.
  \item A model without bias is fitted by least squares to inputs measured
        in degrees Celsius, then to the same inputs in kelvins. Explain
        why the two fitted models can make different predictions, and why
        this cannot happen with a bias.
\end{enumerate}
\solutionlink{affine-spaces:exr:6}
\end{exercise}

\begin{exercise}\label{affine-spaces:exr:7}
On the circle $\mathcal C = \R/\Z$, whose points are the classes
$[x] = x + \Z$, let $v\in\R$ act by $[x] + v = [x + v]$.
\begin{enumerate}[nosep]
  \item Show that this is well defined and satisfies axioms 1 and 2 of
        \cref{affine-spaces:def:affine}, but not axiom 3.
  \item Show that the ``midpoint'' $[(x + y)/2]$ of $[x]$ and $[y]$ depends
        on the representatives $x$ and $y$, and takes exactly two values.
\end{enumerate}
\solutionlink{affine-spaces:exr:7}
\end{exercise}

\begin{exercise}\label{affine-spaces:exr:8}
In homogeneous coordinates
(\cref{affine-spaces:prop:homogeneous}), let $M$ be the matrix of the
affine map $y = Ax + \beta$, with $A$ invertible.
\begin{enumerate}[nosep]
  \item Compute $M^{-1}$ and check that it is the matrix of the inverse
        affine map.
  \item Write the matrix $T_\omega$ of the translation by $\omega$ and
        verify that $MT_\omega$ is the matrix of the same map in the
        coordinates $x' = x - \omega$, with the change of constant term
        $\beta\mapsto\beta + A\omega$ of \cref{affine-spaces:sec:maps}.
\end{enumerate}
\solutionlink{affine-spaces:exr:8}
\end{exercise}

\begin{exercise}\label{affine-spaces:exr:9}
Let $g:\Aff\to\Bff$ be affine, $C\subset\Bff$ convex and $f:C\to\R$ convex.
Show that $g^{-1}(C)$ is convex and that $f\circ g$ is convex on it. Deduce
that for a matrix $X$, a vector $y$ and a convex function
$\varphi:\R^m\to\R$, the function $\theta\mapsto\varphi(X\theta - y)$ is
convex on $\R^n$.
\solutionlink{affine-spaces:exr:9}
\end{exercise}
\end{notechapter}
